\vspace{-1mm}
\section{Related Works} \label{sec:related}


\noindent
\textbf{Invariant Generation and Enforcement.}
Prior research generates and enforces invariants for contract security.
Cider~\cite{liu2022learning} derives arithmetic overflow invariants 
via deep reinforcement learning, 
while InvCon~\cite{liu2022invcon} and InvCon+~\cite{liu2024automated} 
combine dynamic inference with static verification for function-level invariants. 
Trace2Inv~\cite{chen2024demystifying} learns invariants from transaction 
history to prevent attacks. 
Unlike these approaches, which focus on individual 
contracts or specific bugs, {\name} is a general-purpose 
framework that
targets control 
flows across multiple protocol contracts with a one-time 
pre-deployment configuration.




% A significant body of research has focused on generating and enforcing
% invariants for smart contracts to ensure their security.
% % Solythesis~\cite{Solythesis} enforces human defined invariants with its
% % source-to-source compiler. 
% Cider~\cite{liu2022learning} employs deep
% reinforcement learning to derive invariants that prevent arithmetic overflows
% from contract source code. InvCon~\cite{liu2022invcon} and
% InvCon+~\cite{liu2024automated} combine dynamic inference with static
% verification to produce verified, expressive contract invariants based on
% function preconditions and postconditions. 
% %\citeauthor{callens2024temporarily} has designed strategies to prevent the same function from being called twice within a single transaction. 
% Furthermore, Trace2Inv~\cite{chen2024demystifying}
% dynamically learns invariants from transaction history, analyzing the
% effectiveness of $23$ invariants in preventing attacks. These studies
% predominantly focus on invariants tailored to individual contracts or specific
% kinds of vulnerabilities. In contrast, {\name} focuses on control flows across
% multiple contracts within the protocol. Unlike these prior studies, {\name} can
% achieve better true positive and false positive rates with a one-time configuration 
% pre-deployment.
% % without human
% % intervention or the availability of history traces.


% \noindent
% \textbf{Control Flow Restriction.}
% In industry, SphereX~\cite{spherex}, a smart contract security firm, offers
% services to manually restrict control flows within smart contracts. While their
% goals align closely with ours, their approach requires protocol developers to
% manually select which control flows to whitelist or blacklist. This manual
% intervention contrasts sharply with our methodology, which not only automates
% the whitelisting of unseen control flows but also simplifies the overall
% control flow structure, significantly reducing the burden on developers and
% increasing the system's adaptability.


\noindent
\textbf{Control Flow Restriction.}
SphereX~\cite{spherex} in industry offers services to manually restrict control flows, 
requiring developers to explicitly whitelist or 
blacklist paths. 
In contrast, {\name} automates the whitelisting of unseen control 
flows and simplifies the overall structure, 
significantly reducing developer burden and increasing 
system adaptability compared to manual intervention.




\noindent
\textbf{Re-entrancy Attack Defense and Secure Type System.}
Numerous tools restrict control flows to combat re-entrancy attacks. 
Static analysis tools~\cite{luu2016making,torres2018osiris,liu2018reguard,feist2019slither,
zhang2019mpro,bose2022sailfish,ma2021pluto,zheng2022park,wang2024efficiently,
albert2020taming} identify re-entrancy vulnerabilities and apply guards. 
Runtime frameworks like Sereum~\cite{rodler2018sereum} and 
\citeauthor{grossman2017online} protect deployed contracts, 
while \citeauthor{callens2024temporarily} prevent duplicate 
function calls within transactions. Unlike these 
re-entrancy-focused approaches, {\name} adopts broader 
control flow restriction targeting comprehensive vulnerability classes.
\citeauthor{cecchetti2020securing} propose a security type system to 
enforce information flow and re-entrancy controls~\cite{cecchetti2020securing,
cecchetti2021compositional}. Conversely, 
{\name} is a purely runtime system built entirely on the EVM, 
operating without supplementary type systems or language modifications.




% Many researchers have proposed approaches to restrict smart contract control
% flows specifically to combat re-entrancy attacks. Tools like
% Oyente~\cite{luu2016making}, Osiris~\cite{torres2018osiris},
% Reguard~\cite{liu2018reguard}, Slither~\cite{feist2019slither},
% MPro~\cite{zhang2019mpro}, Sailfish~\cite{bose2022sailfish},
% Pluto~\cite{ma2021pluto}, Park~\cite{zheng2022park},
% SliSE~\cite{wang2024efficiently}, \citeauthor{albert2020taming}, 
% employ static analysis and symbolic execution to identify potential re-entrancy
% vulnerabilities, then apply re-entrancy guards as preventive measures. Notably,
% Sereum~\cite{rodler2018sereum} and \citeauthor{grossman2017online} offer
% runtime validation frameworks to protect deployed contracts against re-entrancy
% attacks. \citeauthor{callens2024temporarily} have designed strategies to prevent
% the same function from being called twice within a transaction. Unlike
% these approaches, our work adopts a broader perspective on control flow
% restriction, targeting a more comprehensive set of vulnerabilities beyond just
% re-entrancy.


% \noindent
% \textbf{Secure Type System.}
% \citeauthor{cecchetti2020securing} introduce a security type system that,
% alongside runtime mechanisms, enforces secure information flow and
% re-entrancy controls, even amidst unknown code~\cite{cecchetti2020securing,
% cecchetti2021compositional}. In contrast, our system is designed purely for
% runtime operation built entirely on the EVM, operating without the need for any
% supplementary type system.


